\section{Threats to Validity}
\label{sec:threats}

\paragraph{Construct validity.}
Our goodput SLOs (2\,s/200\,ms and 4\,s/500\,ms) are representative of interactive chat but not
of every application; a tighter TPOT target favours the engine with the fastest decode, and we
report that variant in the supplementary material. Engine-reported metrics have different
semantics across engines (\cref{sec:observability}); we use them only to explain client-side results.
Prefix-cache hit rates are measured per server, not per request.

\paragraph{Internal validity.}
Configuration options are mapped to each engine on a best-effort basis; ``the same'' nominal
setting (e.g., chunk size, sequence limit, n-gram window) reflects each engine's own
implementation. Two configuration defects were found and fixed before the corresponding sweeps
(\cref{sec:engines}); others may remain in less-exercised code paths. Enabling n-gram speculation
on \trt required disabling its overlap scheduler, which confounds that engine's speculative-decoding
result \todo{add a baseline with the overlap scheduler disabled}. \trt's KV-cache and queue
statistics were not observable live, so its scheduler behaviour is inferred from client-side
measurements. \todo{confirm its capacity-scheduler policy from the launch log}

Prompts are not chat-templated for Llama-3.1-8B and Mistral-NeMo-12B. Llama-3.1-8B consequently
generates about six times longer reasoning outputs than the other models, raising its KV pressure
on \texttt{mixed\_real}, and Mistral-NeMo-12B returns zero-token replies for 7--17\% of requests on
every engine. Both effects are identical across engines and therefore do not bias engine
comparisons within a model, but they limit comparisons across models.

\paragraph{External validity.}
Results come from single-GPU deployments of 8--14B-parameter models on two platforms, with fp16
weights and no tensor parallelism. Larger models, quantisation, multi-GPU deployments, and
disaggregated serving may change both absolute numbers and rankings. The \rqref{4} platform also
uses newer engine versions, so hardware and software effects are not separated there. Arrival
processes are synthetic, and real traffic may differ in burstiness and length mix.

\paragraph{Conclusion validity.}
Each configuration was run once, so we report no confidence intervals. We mitigate this by
relying on effects that are large relative to run-to-run variation (often an order of magnitude),
by observing them consistently across arrival rates and models, and by reproducing the main
ranking on a second platform. \todo{add repeated runs for a subset of configurations}
